\chapter{多元函数微分、隐函数与极值}
\section{极限、偏导与全微分}
\ruleitem{多元极限}{$f(x,y)\to L$ 要求从定义域内任意路径趋近同一点都得到 $L$。两条路径不同即可否定；所有直线路径相同仍不足以证明。证明常用 $r=\sqrt{x^2+y^2}$ 控制 $|f-L|\le\omega(r)\to0$。}
\ruleitem{可微定义}{$f(a+h,b+k)=f(a,b)+Ah+Bk+o(\sqrt{h^2+k^2})$。可微推出连续且偏导存在，$A=f_x,B=f_y$；偏导存在未必连续、更未必可微。若偏导在邻域存在且在该点连续，则可微。}
\ruleitem{链式法则}{若 $z=f(u(x,y),v(x,y))$，各映射可微，则 $z_x=f_u u_x+f_v v_x$，$z_y=f_u u_y+f_v v_y$。二阶求导必须同时求“外层偏导”和“内层系数”的导数。若 $z=f(x(t),y(t))$，$z'=\grad f\cdot\bm r'$。}
\ruleitem{混合偏导}{在邻域内相关二阶偏导连续时，$f_{xy}=f_{yx}$。仅两个混合偏导在一点存在，不足以交换。写 $f_{uv}$ 时 $u,v$ 是外层自变量，不是原来的 $x,y$。}
\ruleitem{全微分}{$\dd z=f_x\dd x+f_y\dd y$；$\dd^2z=f_{xx}\dd x^2+2f_{xy}\dd x\dd y+f_{yy}\dd y^2$，此处把 $x,y$ 视为独立变量。复合路径的二阶导数另有 $f_xx''+f_yy''$ 项，不应遗漏。}
\ruleitem{隐函数}{$F(x,y,z)=0$，$F\in C^1$、$F_z\ne0$，则局部
\[
z_x=-F_x/F_z,\quad z_y=-F_y/F_z,\quad
z_{xy}=-\frac{F_{xy}+F_{xz}z_y+F_{yz}z_x+F_{zz}z_xz_y}{F_z}
\]
（后式要求 $C^2$）。方程组隐函数用关于待解变量的雅可比行列式非零作为条件，再解线性方程组。}

\section{方向导数、泰勒与条件极值}
\ruleitem{方向导数}{沿\textbf{单位向量} $\bm e$，可微时 $D_{\bm e}f=\grad f\cdot\bm e$。最大值 $|\grad f|$，方向为梯度方向（梯度非零时）；最小值为其负值。任意方向导数都存在也不能反推可微。}
\ruleitem{齐次函数 Euler 公式}{若 $f(tx,ty)=t^mf(x,y)$（正的 $t$ 在相应域内），且 $f$ 可微，则 $xf_x+yf_y=mf$。三元同理。先检查缩放性质，不可仅凭式子看起来“次数相同”就使用。}
\ruleitem{二元 Taylor}{在 $C^2$ 条件下，$\bm h=(h,k)^T$，
\[
f(\bm a+\bm h)=f(\bm a)+\grad f(\bm a)\cdot\bm h+\frac12\bm h^TH_f(\bm a)\bm h+o(|\bm h|^2).
\]
二次型或 Hessian 正定/负定决定驻点严格极小/极大；不定则为鞍点。}
驻点处令 $A=f_{xx},B=f_{xy},C=f_{yy}$，$\Delta=AC-B^2$：$\Delta>0,A>0$ 为极小；$\Delta>0,A<0$ 为极大；$\Delta<0$ 为鞍点；$\Delta=0$ \textbf{不能判断}。示范：$x^4+y^4$ 极小，$-x^4-y^4$ 极大，$x^4-y^4$ 鞍点，三者二阶量在原点全为零。

\ruleitem{Lagrange 乘子}{约束 $g=0$ 上求 $f$ 极值，正则点 $\grad g\ne0$ 满足 $\grad f=\lambda\grad g$。多约束要求约束梯度线性无关。它给候选点，须比较函数值；约束奇点与边界也需检查。闭有界区域上连续函数必取最值，可按“内部驻点、边界、角点”分别求。}
\ruleitem{几何}{图像曲面切平面 $z-z_0=f_x(x-x_0)+f_y(y-y_0)$；等值面法向为梯度。曲线 $\bm r(t)$ 的切向为 $\bm r'(t)$，正则点的法平面以该切向作法向。}

\section{二阶链式求导真题}
\problem{52}{径向复合函数的 Laplace 算子}{2010 年第 2 届，非数学类，一(4)，\src{02}}
$r=\sqrt{x^2+y^2}>0$，$g(x,y)=f(1/r)$，$f\in C^2$，求 $g_{xx}+g_{yy}$。
\sol 对径向函数 $h(r)$，$r_x=x/r,r_y=y/r$，故 $g_{xx}+g_{yy}=h''+h'/r$。这里
\[h'=-r^{-2}f'(1/r),\quad h''=2r^{-3}f'(1/r)+r^{-4}f''(1/r).\]
从而 $\ans{g_{xx}+g_{yy}=r^{-4}f''(1/r)+r^{-3}f'(1/r)}$。径向 Laplace 公式取决于空间维数；三维的系数为 $2h'/r$，不能混用。

\problem{53}{指数因子消去一阶项}{2012 年第 4 届，非数学类，一(3)，\src{04}}
令 $z=u(x,y)\ee^{ax+by}$，其中 $u_{xy}=0$。求使任意此类 $u$ 均满足 $z_{xy}-z_x-z_y+z=0$ 的 $a,b$。
\sol 提取指数，左端化为 $u_{xy}+(b-1)u_x+(a-1)u_y+(ab-a-b+1)u$。要对任意 $u$ 成立，需 $a=b=1$，常数项也随之为 0。故 $\ans{a=b=1}$。

\problem{54}{隐函数中带交叉分母的组合导数}{2015 年第 7 届，非数学类，一(2)，\src{07}}
$F(x+z/y,y+z/x)=0$ 确定 $z(x,y)$，$xy\ne0$，且 $xF_u+yF_v\ne0$。求 $xz_x+yz_y$（不含 $F$ 及其偏导）。
\sol 分别对 $x,y$ 求导，把含 $z_x,z_y$ 的项集中：
\[
z_x=\frac{y(zF_v-x^2F_u)}{xF_u+yF_v},\qquad
z_y=\frac{x(zF_u-y^2F_v)}{xF_u+yF_v}.
\]
分别乘以 $x,y$ 后相加，分子为 $(z-xy)(xF_u+yF_v)$，所以 $\ans{xz_x+yz_y=z-xy}$。求单个偏导时保留公共分母，目标组合往往会整体约掉。

\problem{55}{两层复合的混合二阶导数}{2023 年第 15 届，非数学 A/B 类，一(2)，共题，\src{15A}\src{15B}}
$z=f(x^2-y^2,xy)$，$f\in C^2$，求 $z_{xy}$。
\sol 令 $u=x^2-y^2,v=xy$，先得 $z_x=2xf_u+yf_v$。对 $y$ 求导：
\[z_{xy}=2x(-2yf_{uu}+xf_{uv})+f_v+y(-2yf_{vu}+xf_{vv}),\]
故 $\ans{z_{xy}=-4xyf_{uu}+2(x^2-y^2)f_{uv}+xyf_{vv}+f_v}$，所有外层偏导均在 $(u,v)$ 处取值。最后的 $f_v$ 来自系数 $y$ 的求导。

\problem{56}{指数内层也要二次求导}{2024 年第 16 届，非数学 A 类一(3)／B 类一(4)，共题，\src{16A}\src{16B}}
$z=f(xy,\ee^{x+y})$，$f\in C^2$，求 $z_{xy}$。
\sol 令 $u=xy,v=\ee^{x+y}$，$z_x=yf_u+vf_v$。再对 $y$ 求导并合并：
\[
\ans{z_{xy}=xyf_{uu}+(x+y)vf_{uv}+v^2f_{vv}+f_u+vf_v},\quad v=\ee^{x+y}.
\]
其中 $f_u$ 来自 $y$ 的导数，$vf_v$ 来自 $v$ 的导数；它们是二阶求导最常漏的项。

\problem{57}{先求隐函数值，再求偏导}{2025 年第 17 届，非数学 A/B 类，一(4)，共题，\src{17A}\src{17B}}
$xz=f(2x-y,xyz)$，$f(2,0)=3,f_u(2,0)=5$，$f\in C^1$。求 $z_x(1,0)$。
\sol 先由原式得 $z(1,0)=3$。固定 $y$ 求 $x$ 偏导：$z+xz_x=2f_u+f_v(yz+xyz_x)$。在 $(1,0)$，所有带 $y$ 的项消失，故 $3+z_x=10$，$\ans{z_x(1,0)=7}$。无需知道 $f_v(2,0)$。

\section{梯度、极值与约束真题}
\problem{58}{双约束最值可先消元}{2024 年第 16 届，非数学 B 类，二，\src{16B}}
在 $x-y+z=1$、$x^2+y^2=1$ 下，求 $x+2y+3z$ 的最大值及取值点。
\sol 消去 $z=1-x+y$，目标为 $3-2x+5y$。Cauchy 不等式给出 $-2x+5y\le\sqrt{29}\sqrt{x^2+y^2}=\sqrt{29}$，等号在 $(x,y)=(-2,5)/\sqrt{29}$。
故最大值 $\ans{3+\sqrt{29}}$，取值点 $\ans{(-2/\sqrt{29},5/\sqrt{29},1+7/\sqrt{29})}$。若用乘子法，消元后的方程是 $(-2,5)=2\lambda(x,y)$，同样得到此点。

\problem{59}{方向导数信息确定曲面，再求面积}{2025 年第 17 届，非数学 B 类，二，\src{17B}}
$f(x,y)=ax^2+by^2+1$。在 $(4,3)$ 处沿指向原点的方向取得最大方向导数 10。求 $a,b$，并求曲面 $z=f(x,y)$ 位于 $z=0$ 上方部分的面积。
\sol 指向原点的单位方向为 $(-4/5,-3/5)$，最大值为梯度模长，所以 $\grad f(4,3)=10(-4/5,-3/5)=(-8,-6)$。而梯度是 $(8a,6b)$，得 $\ans{a=b=-1}$。
曲面 $z=1-x^2-y^2$，投影为单位圆盘，面积
\[S=\iint_{x^2+y^2\le1}\sqrt{1+4x^2+4y^2}\dd x\dd y
=2\pi\int_0^1r\sqrt{1+4r^2}\dd r=\ans{\frac\pi6(5\sqrt5-1)}.\]
用方向向量前必须单位化，直接把 $(-4,-3)$ 当方向余弦会把梯度放大 5 倍。

\problem{60}{隐函数两个分支的局部极值}{2025 年第 17 届，非数学 B 类，三，\src{17B}}
由 $2x^2+2y^2+z^2+8xz-z+8=0$ 确定的局部隐函数 $z(x,y)$，求极值。
\sol 记左端为 $F$。驻点要求 $F_x=4x+8z=0$、$F_y=4y=0$，所以 $x=-2z,y=0$。代回得 $7z^2+z-8=0$，即 $z=1,-8/7$。
在 $(-2,0,1)$，$F_z=2z+8x-1=-15\ne0$；驻点处 $z_{xx}=z_{yy}=-4/F_z=4/15$、$z_{xy}=0$，故为极小值 $\ans1$。
在 $(16/7,0,-8/7)$，$F_z=15$，Hessian 负定，故为极大值 $\ans{-8/7}$。两者属于不同局部支，不能因“极小值比另一个极大值大”而认为矛盾。

\problem{61}{沿线段把多元问题化为一元}{2018 年第 10 届，非数学类，五，\src{10}}
$f$ 在域内可微，$f_x^2+f_y^2\le M$。若线段 $AB$ 全在域内，证明 $|f(B)-f(A)|\le\sqrt M\,|AB|$。
\sol 令 $\phi(t)=f(A+t(B-A))$，$0\le t\le1$。链式法则及 Cauchy 给出 $|\phi'(t)|=|\grad f\cdot(B-A)|\le\sqrt M|AB|$。用 Lagrange 或积分即得。
\trap{线段必须留在定义域中。任意连通域不必包含连接两点的直线段；此时只能沿域内路径估计长度。}
